

\bta{2006乙B}

\begin{enumerate}
	%\renewcommand{\labelenumi}{\arabic{enumi}.}
	% A(\Alph) a(\alph) I(\Roman) i(\roman) 1(\arabic)
	%设定全局标号series=example	%引用全局变量resume=example
	%[topsep=-0.3em,parsep=-0.3em,itemsep=-0.3em,partopsep=-0.3em]
	%可使用leftmargin调整列表环境左边的空白长度 [leftmargin=0em]
	\item
（$30'$）粒子以能量$E$入射方势垒，$V(x)=\begin{cases}V_0>0,&0\le x\le a\\0,&x<0\ \text{或}\ x>a\end{cases}$。设能量$E<V_0$，求透射系数$T$。


\banswer{
假定 $0<E<V_0$。令 $k=\sqrt{2mE}/\hbar$、$\kappa=\sqrt{2m(V_0-E)}/\hbar$，取
\[
 \psi_{\rm I}=e^{ikx}+r e^{-ikx},\quad
 \psi_{\rm II}=Ae^{\kappa x}+Be^{-\kappa x},\quad
 \psi_{\rm III}=t e^{ikx}.
\]
界面连续条件为
\[
 \begin{cases}
 1+r=A+B,\\ ik(1-r)=\kappa(A-B),\\
 Ae^{\kappa a}+Be^{-\kappa a}=t e^{ika},\\
 \kappa(Ae^{\kappa a}-Be^{-\kappa a})=ik t e^{ika}.
 \end{cases}
\]
先由后两式解出 $Ae^{\kappa a}=t e^{ika}(1+ik/\kappa)/2$、
$Be^{-\kappa a}=t e^{ika}(1-ik/\kappa)/2$，代回前两式消去 $r$，得
\[
 t=\frac{e^{-ika}}{\cosh(\kappa a)
                  +i\dfrac{\kappa^2-k^2}{2k\kappa}\sinh(\kappa a)}.
\]
左右两侧相同，故 $T=|t|^2$。使用 $\cosh^2u=1+\sinh^2u$，
\[
 T^{-1}=1+\frac{(k^2+\kappa^2)^2}{4k^2\kappa^2}\sinh^2(\kappa a)
       =1+\frac{V_0^2}{4E(V_0-E)}\sinh^2(\kappa a).
\]
所以
\[
 \boxed{T=\left[1+\frac{V_0^2}{4E(V_0-E)}\sinh^2(\kappa a)\right]^{-1}}.
\]
在厚势垒极限 $\kappa a\gg1$，$\sinh^2(\kappa a)\simeq e^{2\kappa a}/4$，
得到 $T\simeq16E(V_0-E)V_0^{-2}e^{-2\kappa a}$。
}


\item 
（$30'$）粒子在一维对称无限深方势阱$(-\frac{a}{2}\le x\le \frac{a}{2})$中运动。设$t=0$时粒子所处状态为$\psi(x,t=0)=\frac{1}{\sqrt2}[\varphi_1(x)+\varphi_2(x)]$，其中$\varphi_n(x)$为系统第$n$个能量本征态。求$t>0$时的以下量：
\begin{enumerate}
	%\renewcommand{\labelenumi}{\arabic{enumi}.}
	% A(\Alph) a(\alph) I(\Roman) i(\roman) 1(\arabic)
	%设定全局标号series=example	%引用全局变量resume=example
	%[topsep=-0.3em,parsep=-0.3em,itemsep=-0.3em,partopsep=-0.3em]
	%可使用leftmargin调整列表环境左边的空白长度 [leftmargin=0em]
	\item
概率密度$|\psi(x,t)|^2$；

\item 
能量的可取值及相应的概率。
	
\end{enumerate}


\banswer{
	取
	\[
		\phi_1(x)=\sqrt{\frac2a}\cos\frac{\pi x}{a},
		\qquad
		\phi_2(x)=\sqrt{\frac2a}\sin\frac{2\pi x}{a},
		\qquad
		E_n=\frac{n^2\pi^2\hbar^2}{2ma^2}.
	\]
	则
	\[
		\psi(x,t)=\frac1{\sqrt2}
		\left[e^{-iE_1t/\hbar}\phi_1(x)+e^{-iE_2t/\hbar}\phi_2(x)\right].
	\]
	\begin{enumerate}
		\item 概率密度为
		\[
			\boxed{|\psi(x,t)|^2=\frac12\left[\phi_1^2(x)+\phi_2^2(x)
			+2\phi_1(x)\phi_2(x)\cos\frac{(E_2-E_1)t}{\hbar}\right]}.
		\]
		\item 能量可取 $E_1=\pi^2\hbar^2/(2ma^2)$ 和 $E_2=4E_1$，相应概率均为 $1/2$，且不随时间改变。
	\end{enumerate}
}


\item 
（$30'$）设氢原子所处状态为$\psi(r,\theta,\phi,S_z)=\begin{pmatrix}\frac{1}{2}R_{21}(r)Y_{11}(\theta,\phi)\\-\frac{\sqrt3}{2}R_{21}(r)Y_{10}(\theta,\phi)\end{pmatrix}$。
\begin{enumerate}
	%\renewcommand{\labelenumi}{\arabic{enumi}.}
	% A(\Alph) a(\alph) I(\Roman) i(\roman) 1(\arabic)
	%设定全局标号series=example	%引用全局变量resume=example
	%[topsep=-0.3em,parsep=-0.3em,itemsep=-0.3em,partopsep=-0.3em]
	%可使用leftmargin调整列表环境左边的空白长度 [leftmargin=0em]
	\item
求轨道角动量$z$分量$\hat{L}_z$和自旋角动量$z$分量$\hat{S}_z$的平均值；

\item 
求总磁矩$\vec{M}=-\frac{e}{2\mu}\vec{L}-\frac{e}{\mu}\vec{S}$的$z$分量的平均值。
	
\end{enumerate}

\banswer{
	把题给两分量分别理解为 $S_z=+\hbar/2$ 与 $S_z=-\hbar/2$ 的振幅。两分量概率为 $1/4$ 与 $3/4$。
	\begin{enumerate}
		\item 因 $Y_{11}$、$Y_{10}$ 的 $L_z$ 本征值分别为 $\hbar$、$0$，
		\[
			\boxed{\langle L_z\rangle=\frac14\hbar}.
		\]
		同时
		\[
			\boxed{\langle S_z\rangle=\frac14\frac\hbar2
			+\frac34\left(-\frac\hbar2\right)=-\frac\hbar4}.
		\]

		\item
		\begin{align*}
			\langle M_z\rangle
			&=-\frac{e}{2\mu}\langle L_z\rangle
			-\frac{e}{\mu}\langle S_z\rangle
			=\boxed{\frac{e\hbar}{8\mu}}.
		\end{align*}
	\end{enumerate}
}

\newpage
\item 
对于一维谐振子的基态，求坐标和动量的不确定度的乘积$\Delta x\cdot\Delta p$。

\banswer{
	谐振子基态由 $a|0\rangle=0$ 决定。写
	\[
	 x=\sqrt{\frac{\hbar}{2m\omega}}(a+a^\dagger),\qquad
	 p=i\sqrt{\frac{m\hbar\omega}{2}}(a^\dagger-a).
	\]
	基态中 $\langle a\rangle=\langle a^\dagger\rangle=0$，
	$\langle aa^\dagger\rangle=1$，而 $\langle a^\dagger a\rangle=
	\langle a^2\rangle=\langle a^{\dagger2}\rangle=0$。展开平方得到
	\[
		\langle x\rangle=\langle p\rangle=0,\qquad
		\langle x^2\rangle=\frac{\hbar}{2m\omega},\qquad
		\langle p^2\rangle=\frac{m\hbar\omega}{2}.
	\]
	因此
	\[
		\boxed{\Delta x=\sqrt{\frac{\hbar}{2m\omega}},
		\qquad\Delta p=\sqrt{\frac{m\hbar\omega}{2}},
		\qquad\Delta x\,\Delta p=\frac\hbar2}.
	\]
}


\item 
（$30'$）两个自旋为$\frac{1}{2}$非全同粒子，自旋间相互作用为$\hat{H}=J\vec{s}_1\cdot\vec{s}_2$，其中$\vec{s}_1$和$\vec{s}_2$分别为粒子1和粒子2的自旋算符。设$t=0$时粒子1的自旋沿$z$轴正方向，粒子2的自旋沿$z$轴负方向。求$t>0$时，测到粒子2的自旋仍处于$z$轴负方向的概率。


\banswer{
	初态可按总自旋单态、三重态展开：
	\[
		|\uparrow\downarrow\rangle
		=\frac1{\sqrt2}\left(|1,0\rangle+|0,0\rangle\right).
	\]
	由
	\[
		\vec s_1\cdot\vec s_2
		=\frac12(S^2-s_1^2-s_2^2),
	\]
	三重态与单态能量分别为 $E_T=J\hbar^2/4$、$E_S=-3J\hbar^2/4$。所以
	\[
		|\psi(t)\rangle=\frac1{\sqrt2}
		\left(e^{-iE_Tt/\hbar}|1,0\rangle
		+e^{-iE_St/\hbar}|0,0\rangle\right).
	\]
	投影回 $|\uparrow\downarrow\rangle$，即粒子 $2$ 自旋仍向下的概率为
	\[
		\boxed{P_{2\downarrow}(t)=
		\frac14\left|e^{-iE_Tt/\hbar}+e^{-iE_St/\hbar}\right|^2
		=\cos^2\left(\frac{J\hbar t}{2}\right)}.
	\]
}


	
\end{enumerate}
